\section{Morphisms of sites}
\label{section-morphism-sites}

\begin{definition}
\label{definition-morphism-sites}
Let $\mathcal{C}$ and $\mathcal{D}$ be sites.
A {\it morphism of sites} $f : \mathcal{D} \to \mathcal{C}$
is given by a continuous functor $u : \mathcal{C} \to \mathcal{D}$
such that the functor $u_s$ is exact.
\end{definition}

\noindent
Notice how the functor $u$ goes in the direction {\it opposite}
the morphism $f$. If $f \leftrightarrow u$ is a morphism of sites
then we use the notation $f^{-1} = u_s$ and $f_* = u^s$.
The functor $f^{-1}$ is called the {\it pullback functor} and
the functor $f_*$ is called the {\it pushforward functor}.
As in topology we have the following adjointness property
$$
\Mor_{\Sh(\mathcal{D})}(f^{-1}\mathcal{G}, \mathcal{F})
=
\Mor_{\Sh(\mathcal{C})}(\mathcal{G}, f_*\mathcal{F})
$$
The motivation for this definition comes from the following
example.

\begin{example}
\label{example-continuous-map}
Let $f : X  \to Y$ be a continuous map of topological spaces.
Recall that we have sites $X_{Zar}$ and $Y_{Zar}$,
see Example \ref{example-site-topological}. Consider the functor
$u : Y_{Zar} \to X_{Zar}$, $V \mapsto f^{-1}(V)$.
This functor is clearly continuous because inverse images of
open coverings are open coverings. (Actually, this depends on how
you chose sets of coverings for $X_{Zar}$ and $Y_{Zar}$.
But in any case the functor is quasi-continuous, see Remark
\ref{remark-quasi-continuous}.)
It is easy to verify that
the functor $u^s$ equals the usual pushforward functor $f_*$
from topology. Hence, since $u_s$ is an adjoint and since
the usual topological pullback functor $f^{-1}$ is an adjoint as well,
we get a canonical isomorphism $f^{-1} \cong u_s$. Since $f^{-1}$
is exact we deduce that $u_s$ is exact. Hence $u$ defines a morphism
of sites $f : X_{Zar} \to Y_{Zar}$, which we may denote
$f$ as well since we've already seen the functors $u_s, u^s$ agree
with their usual notions anyway.
\end{example}

\begin{example}
\label{example-finer-topology}
Let $\mathcal{C}$ be a category. Let
$$
\text{Cov}(\mathcal{C}) \supset \text{Cov}'(\mathcal{C})
$$
be two sets of families of morphisms with fixed target
which turn $\mathcal{C}$ into a site. Denote $\mathcal{C}_\tau$
the site corresponding to $\text{Cov}(\mathcal{C})$ and
$\mathcal{C}_{\tau'}$ the site corresponding to $\text{Cov}'(\mathcal{C})$.
We claim the identity functor on $\mathcal{C}$ defines a morphism of sites
$$
\epsilon : \mathcal{C}_\tau \longrightarrow \mathcal{C}_{\tau'}
$$
Namely, observe that $\text{id} : \mathcal{C}_{\tau'} \to \mathcal{C}_\tau$
is continuous as every
$\tau'$-covering is a $\tau$-covering. Thus the functor
$\epsilon_* = \text{id}^s$ is the identity functor on underlying
presheaves. Hence the left adjoint $\epsilon^{-1}$
of $\epsilon_*$ sends a $\tau'$-sheaf $\mathcal{F}$
to the $\tau$-sheafification of $\mathcal{F}$ (by the
universal property of sheafification).
Finite limits of $\tau'$-sheaves agree with finite limits
of presheaves (Lemma \ref{lemma-limit-sheaf}) and $\tau$-sheafification
commutes with finite limits (Lemma \ref{lemma-sheafification-exact}).
Thus $\epsilon^{-1}$ is left exact. Since $\epsilon^{-1}$
is a left adjoint it is also right exact
(Categories, Lemma \ref{categories-lemma-exact-adjoint}).
Thus $\epsilon^{-1}$ is exact and we have checked all the
conditions of Definition \ref{definition-morphism-sites}.
\end{example}

\begin{lemma}
\label{lemma-composition-morphisms-sites}
\begin{slogan}
Composition of site functors respects continuity of site functors.
\end{slogan}
Let $\mathcal{C}_i$, $i = 1, 2, 3$ be sites. Let
$u : \mathcal{C}_2 \to \mathcal{C}_1$ and
$v : \mathcal{C}_3 \to \mathcal{C}_2$ be continuous functors
which induce morphisms of sites. Then the functor
$u \circ v : \mathcal{C}_3 \to \mathcal{C}_1$ is continuous and
defines a morphism of sites $\mathcal{C}_1 \to \mathcal{C}_3$.
\end{lemma}

\begin{proof}
It is immediate from the definitions that $u \circ v$ is a continuous functor.
In addition, we clearly have $(u \circ v)^p = v^p \circ u^p$, and hence
$(u \circ v)^s = v^s \circ u^s$. Hence functors $(u \circ v)_s$  and
$u_s \circ v_s$ are both left adjoints of $(u \circ v)^s$. Therefore
$(u \circ v)_s \cong u_s \circ v_s$ and we conclude that $(u \circ v)_s$
is exact as a composition of exact functors.
\end{proof}

\begin{definition}
\label{definition-composition-morphisms-sites}
Let $\mathcal{C}_i$, $i = 1, 2, 3$ be sites. Let
$f : \mathcal{C}_1 \to \mathcal{C}_2$ and
$g : \mathcal{C}_2 \to \mathcal{C}_3$ be morphisms of sites
given by continuous functors $u : \mathcal{C}_2 \to \mathcal{C}_1$
and $v : \mathcal{C}_3 \to \mathcal{C}_2$. The {\it composition}
$g \circ f$ is the morphism of sites corresponding to the
functor $u \circ v$.
\end{definition}

\noindent
In this situation we have $(g \circ f)_* = g_* \circ f_*$ and
$(g \circ f)^{-1} = f^{-1} \circ g^{-1}$ (see proof of
Lemma \ref{lemma-composition-morphisms-sites}).

\begin{lemma}
\label{lemma-directed-morphism}
Let $\mathcal{C}$ and $\mathcal{D}$ be sites. Let
$u : \mathcal{C} \to \mathcal{D}$ be continuous.
Assume all the categories $(\mathcal{I}_V^u)^{opp}$ of
Section \ref{section-functoriality-PSh}
are filtered. Then $u$ defines a morphism of sites $\mathcal{D} \to
\mathcal{C}$, in other words $u_s$ is exact.
\end{lemma}

\begin{proof}
Since $u_s$ is the left adjoint of $u^s$ we see that $u_s$ is right
exact, see Categories, Lemma \ref{categories-lemma-exact-adjoint}.
Hence it suffices to show that $u_s$ is left exact. In other words
we have to show that $u_s$ commutes with finite limits.
Because the categories $\mathcal{I}_Y^{opp}$ are filtered
we see that $u_p$ commutes with finite limits, see
Categories, Lemma \ref{categories-lemma-directed-commutes}
(this also uses the description of limits in $\textit{PSh}$,
see Section \ref{section-limits-colimits-PSh}).
And since sheafification commutes with finite limits as well
(Lemma \ref{lemma-sheafification-exact}) we conclude because
$u_s = \# \circ u_p$.
\end{proof}

\begin{proposition}
\label{proposition-get-morphism}
Let $\mathcal{C}$ and $\mathcal{D}$ be sites. Let
$u : \mathcal{C} \to \mathcal{D}$ be continuous.
Assume furthermore the following:
\begin{enumerate}
\item the category $\mathcal{C}$ has a final object $X$ and
$u(X)$ is a final object of $\mathcal{D}$ , and
\item the category $\mathcal{C}$ has fibre products and
$u$ commutes with them.
\end{enumerate}
Then $u$ defines a morphism of sites $\mathcal{D} \to
\mathcal{C}$, in other words $u_s$ is exact.
\end{proposition}

\begin{proof}
This follows from Lemmas \ref{lemma-directed} and
\ref{lemma-directed-morphism}.
\end{proof}

\begin{remark}
\label{remark-explain-left-exact}
The conditions of Proposition \ref{proposition-get-morphism} above
are equivalent to saying that $u$ is left exact, i.e., commutes
with finite limits. See
Categories, Lemmas
\ref{categories-lemma-finite-limits-exist} and
\ref{categories-lemma-characterize-left-exact}.
It seems more natural to phrase it in terms of final objects
and fibre products since this seems to have more geometric meaning
in the examples.
\end{remark}

\noindent
Lemma \ref{lemma-continuous-with-continuous-left-adjoint} will
provide another way to prove a continuous functor
gives rise to a morphism of sites.

\begin{remark}
\label{remark-quasi-continuous-morphism-sites}
(Skip on first reading.)
Let $\mathcal{C}$ and $\mathcal{D}$ be sites. Analogously to
Definition \ref{definition-morphism-sites} we say that
a {\it quasi-morphism of sites $f : \mathcal{D} \to \mathcal{C}$}
is given by a quasi-continuous functor $u : \mathcal{C} \to \mathcal{D}$
(see Remark \ref{remark-quasi-continuous}) such that $u_s$ is exact.
The analogue of Proposition \ref{proposition-get-morphism} in this
setting is obtained by replacing the word ``continuous''
by the word ``quasi-continuous'', and replacing the word
``morphism'' by ``quasi-morphism''. The proof is literally the
same.
\end{remark}

\noindent
In Definition \ref{definition-morphism-sites} the condition that $u_s$
be exact cannot be omitted. For example, the conclusion of the following
lemma need not hold if one only assumes that $u$ is continuous.

\begin{lemma}
\label{lemma-morphism-of-sites-covering}
Let $f : \mathcal{D} \to \mathcal{C}$ be a morphism of sites given by the
functor $u : \mathcal{C} \to \mathcal{D}$. Given any object $V$ of
$\mathcal{D}$ there exists a covering $\{V_j \to V\}$ such that for every
$j$ there exists a morphism $V_j \to u(U_j)$ for some object $U_j$
of $\mathcal{C}$.
\end{lemma}

\begin{proof}
Since $f^{-1} = u_s$ is exact we have $f^{-1}* = *$ where $*$ denotes the
final object of the category of sheaves
(Example \ref{example-singleton-sheaf}).
Since $f^{-1}* = u_s*$ is the sheafification of $u_p*$ we see
there exists a covering $\{V_j \to V\}$ such that $(u_p*)(V_j)$
is nonempty. Since $(u_p*)(V_j)$ is a colimit over the category
$\mathcal{I}^u_{V_j}$ whose objects are morphisms $V_j \to u(U)$
the lemma follows.
\end{proof}



























